% SGA 6, Exposición IV: fragmento español alineado con la fuente.
% Parte 1 de 7. Alcance: apertura, sección 1, líneas 6118--6208 inclusive.
% Fuente francesa: sga6_fr_r1/src/fr.tex.
% SHA-256 de la fuente francesa completa:
% 0D2158B9FE3B466309A92165727A0B191422E05FBAA0DACD5806E8322C54BA1A.
\providecommand{\sisi}[2]{#2}

\section*{Exposición IV. Grupos de Grothendieck de los topoi anillados}
\begin{center}
por L. Illusie
\end{center}

\section*{1. Recordatorios y generalidades sobre los grupos de Grothendieck}

\textbf{1.1.} Sea \(C\) una categoría triangulada. Recordemos (I 6.3 y SGA 5 VIII 2) que una aplicación \(f\) de \(\Ob C\) en un grupo abeliano \(G\) se dice \emph{aditiva} si se tiene \(f(E)=f(E')+f(E'')\) para todo triángulo distinguido \(E'\longrightarrow E\longrightarrow E''\longrightarrow E'[1]\). Las aplicaciones aditivas de \(\Ob C\) en \(G\) forman un grupo abeliano, que depende funtorialmente de \(G\). El funtor así obtenido está representado por un grupo abeliano \(k(C)\) y una aplicación aditiva universal \(\cl_C:\Ob C\longrightarrow k(C)\) (denotada \(\cl\) cuando no hay riesgo de confusión).

El grupo \(k(C)\) depende funtorialmente de \(C\) con respecto a los funtores exactos. Si \(C,C'\) son categorías trianguladas, dos funtores exactos isomorfos de \(C\) en \(C'\) inducen el mismo homomorfismo \(k(C)\to k(C')\); en particular, si \(u:C\to C'\) es un funtor exacto que es una equivalencia de categorías, \(k(u):k(C)\to k(C')\) es un isomorfismo.

\textbf{Lema 1.2.} Sean \(C\) una categoría triangulada, y \(L\in\Ob C\).
\begin{enumerate}[label=\alph*)]
\item Para todo \(n\in\mathbb Z\), se tiene
\[
 \cl(L[n])=(-1)^n\cl(L).
\]
Si \(L',L''\) son objetos de \(C\) tales que \(L\simeq L'\oplus L''\), entonces se tiene \(\cl(L)=\cl(L')+\cl(L'')\).
\item Supongamos que existe \(n\in\mathbb Z\) tal que \(\coprod_{i\geq0} L[2ni]\) sea representable. Entonces se tiene \(\cl(L)=0\).
\end{enumerate}

\emph{Demostración.} La afirmación a) se deduce directamente de las definiciones (cf. SGA 5 VIII 2). Para b), basta observar que se tiene
\[
 \coprod_{i\geq0} L[2ni]\xrightarrow{\sim}
 L\oplus\left(\coprod_{i\geq0}L[2ni]\right)[2n],
\]
y aplicar a).

\textbf{Lema 1.3.} Sean \(C\) una categoría triangulada, \(A\) una subcategoría gruesa \(([\mathrm V]\ \mathrm I\ 2.1)\). Los funtores de inclusión y de paso al cociente definen una sucesión exacta
\[
 k(A)\longrightarrow k(C)\longrightarrow k(C/A)\longrightarrow 0.
\]

\emph{Demostración.} Véase (SGA 5 VIII 3.1).

\textbf{Lema 1.4.} Sean \(A\) una categoría aditiva (resp. abeliana), y \(f\) una función aditiva (1.1) sobre \(\Ob K^b(A)\) (resp. \(\Ob D^b(A)\)). Para \(E\in\Ob K^b(A)\) (resp. \(D^b(A)\)) se tiene
\[
 f(E)=\sum(-1)^i f(E^i)
\]
(resp.
\[
 f(E)=\sum(-1)^i f(H^i(E))
\]).

\emph{Demostración.} Se deja como ejercicio al lector.

\textbf{1.5.} Sean \(C\) una categoría abeliana y \(A\) una subcategoría plena de \(C\) estable por sumas finitas. Una aplicación \(f\) de \(\Ob A\) en un grupo abeliano \(G\) se dice \emph{aditiva} si se tiene \(\sum(-1)^i f(E^i)=0\) para todo complejo acíclico \(E\) de \(C\), de grados acotados y cuyas componentes pertenecen a \(A\). Cuando \(A\) es estable por núcleo de epimorfismo (I 1.2), equivale a decir que se tiene \(f(L)=f(L')+f(L'')\) para toda sucesión exacta \(0\longrightarrow L'\longrightarrow L\longrightarrow L''\longrightarrow0\) de \(C\) cuyos términos pertenecen a \(A\). Las aplicaciones aditivas de \(\Ob A\) en \(G\) forman un grupo abeliano que depende funtorialmente de \(G\). Denotemos por \(K^{b,\emptyset}(A)\) la subcategoría gruesa de \(K^b(A)\) formada por los complejos acíclicos. Se deduce fácilmente de (1.3) y (1.4) que la aplicación
\begin{equation}\tag{1.5.1}
 \cl_A:\Ob A\longrightarrow k\bigl(K^b(A)/K^{b,\emptyset}(A)\bigr)
\end{equation}
inducida por \(\cl_{K^b(A)/K^{b,\emptyset}(A)}\) (1.1) es una aplicación aditiva universal, es decir, que el par \(\left(k\bigl(K^b(A)/K^{b,\emptyset}(A)\bigr),\cl_A\right)\) representa el funtor que asocia a cada grupo abeliano \(G\) el grupo abeliano de las aplicaciones aditivas de \(\Ob A\) en \(G\). Pondremos
\begin{equation}\tag{1.5.2}
 k\bigl(K^b(A)/K^{b,\emptyset}(A)\bigr)=k(A)
\end{equation}
(se escribirá \(k(A,C)\) en lugar de \(k(A)\) cuando se quiera evitar toda ambigüedad sobre la categoría abeliana ambiente)\footnote{La notación \(k(A)\) se había introducido en (I 6.3) para designar el grupo \(k(K^b(A))\), que no es en general isomorfo al grupo definido por (1.5.2). Cambiamos, pues, aquí de convención: en la presente exposición la notación \(k(A)\) designará exclusivamente el grupo definido por (1.5.2).}.
En particular, para \(A=C\), se tiene
\begin{equation}\tag{1.5.3}
 k(D^b(C))=k(C).
\end{equation}

\textbf{Lema 1.6.} Sean \(C\) una categoría abeliana, \(A\) una subcategoría gruesa (es decir, una subcategoría plena estable por núcleos, conúcleos y extensiones). Denotemos por \(D_A(C)\) la subcategoría plena de \(D(C)\) formada por los objetos \(E\) tales que \(H^i(E)\in\Ob A\) para todo \(i\). Entonces \(D_A(C)\) es una subcategoría triangulada \(([\mathrm V]\ \mathrm I\ 1.2.3)\) de \(D(C)\), y el funtor canónico \(D^b(A)\longrightarrow D_A^b(C)=D^b(C)\cap D_A(C)\) induce un isomorfismo
\[
 k(A)\xrightarrow{\sim} k(D_A^b(C)),
\]
donde \(k(A)\) está definido por (1.5.2).

\emph{Demostración.} La primera afirmación es consecuencia inmediata de la sucesión exacta de cohomología. Para la segunda, consideremos la aplicación \(f:\Ob D_A^b(C)\longrightarrow k(A)\) definida por
\[
 f(E)=\sum(-1)^i\cl_A(H^i(E)),
\]
donde \(\cl_A\) es la aplicación (1.5.1). Según la sucesión exacta de cohomología, \(f\) es aditiva (1.1), por lo que define un homomorfismo
\[
 \varphi:k(D_A^b(C))\longrightarrow k(A).
\]
Se deduce de inmediato de (1.4) que \(\varphi\) es inverso del homomorfismo canónico
\[
 k(A)\longrightarrow k(D_A^b(C)).
\]

\textbf{Observación 1.6.1.} En la situación de (1.6), el funtor canónico \(D^b(A)\longrightarrow D_A^b(C)\) no es en general una equivalencia de categorías, como muestra el ejemplo en que \(S\) es un esquema, \(C\) la categoría de los \(\mathcal O_S\)-módulos, \(A\) la categoría de los \(\mathcal O_S\)-módulos cuasicoherentes (II Ap. I 0.3). Véase, no obstante, el criterio (II 2.1.3).

\textbf{Definición 1.7.} Sean \(C\) una categoría triangulada (resp. abeliana), y \(C'\) una parte de \(\Ob C\). Diremos que \(C'\) es un \emph{subconjunto exacto} (cf. EGA III 3.1.1) si, para todo triángulo distinguido \(E'\to E\to E''\to E'[1]\) (resp. toda sucesión exacta \(0\to E'\to E\to E''\to0\)) tal que dos de los términos \(E',E,E''\) pertenezcan a \(C'\), el tercero también pertenece a dicho subconjunto. Diremos que \(C'\) es \emph{denso} en \(\Ob C\) si \(\Ob C\) es el menor subconjunto exacto de \(\Ob C\) que contiene a \(C'\).

\textbf{Lema 1.8.} Sean \(C\) una categoría triangulada (resp. abeliana) y \(C'\) una parte densa de \(\Ob C\). Entonces los elementos \(\cl_C(E)\), para \(E\in C'\), generan \(k(C)\).

\emph{Demostración.} Pongamos \(C'=C_0\), y, para \(i\geq1\), definamos \(C_i\) por recurrencia como el conjunto de los objetos de \(C\) que figuran en un triángulo distinguido (resp. una sucesión exacta corta) dos de cuyos términos pertenecen a \(C_{i-1}\). Los \(C_i\) forman una filtración creciente de \(\Ob C\), que, por hipótesis, es exhaustiva. Para \(i\geq0\), denotemos por \(k(C)_i\) el subgrupo de \(k(C)\) generado por los elementos \(\cl_C(E)\) para \(E\in C_i\). Es inmediato a partir de las definiciones que se tiene, para todo \(i\geq0\), \(k(C)_i=k(C)_{i+1}\), y por tanto \(k(C)_i=k(C)_0\). Sea entonces \(x\in k(C)\); se puede escribir, como se ve fácilmente, \(x=\cl(E')-\cl(E'')\) para \(E',E''\in\Ob C\) (incluso puede disponerse, en el caso \sisi{no respé}{no resp.}, que \(E''=0\)). Como la filtración \((C_i)\) es exhaustiva, existe un entero \(i\geq0\) tal que \(E'\) y \(E''\) pertenecen a \(C_i\); se tiene entonces \(x\in k(C)_i=k(C)_0\), y se concluye.

\textbf{Ejercicio 1.9.} Sean \(C\) una categoría abeliana, \(A\) una subcategoría plena estable por sumas finitas, subobjetos y cocientes, y tal que \(\Ob A\) sea denso en \(\Ob C\). Entonces el homomorfismo \(k(A,C)\longrightarrow k(C)\) definido por la función aditiva \(E\mapsto\cl_C(E)\) sobre \(\Ob A\) (cf. (1.5)) es un isomorfismo.

\emph{Indicación}: Poner \(A=A_0\), y definir \(A_i\), para \(i\geq1\), como la subcategoría plena de \(C\) formada por los objetos de \(C\) que son extensiones de dos objetos de \(C_{i-1}\). Probar que \(A_i\) satisface las mismas hipótesis que \(A\), y mostrar después que el homomorfismo canónico \(\varinjlim_i k(A_i,C)\longrightarrow k(C)\) es un isomorfismo (la conclusión buscada se deduce de inmediato).
